\clearpage\SourceLeaf{172}\pdfbookmark[0]{A New Method for Finding the Values of Integrals by Approximation}{seventh-paper}\thispagestyle{empty}
\begin{center}
A NEW METHOD

\medskip
FOR FINDING THE VALUES OF INTEGRALS

BY APPROXIMATION

BY

CARL FRIEDRICH GAUSS

PRESENTED TO THE ROYAL SOCIETY OF SCIENCES ON 16 SEPTEMBER 1814.

\par\bigskip\rule{.7\linewidth}{.4pt}\par\medskip
Commentationes societatis regiae scientiarum Gottingensis recentiores. Vol.~\textsc{iii}.

Gottingae \textsc{mdcccxvi}.
\par\medskip\rule{.7\linewidth}{.4pt}
\end{center}
\clearpage\fancyhead[L]{\small APPROXIMATE VALUES OF INTEGRALS}
\SourceLeaf{174}
\begin{center}
A NEW METHOD

FOR FINDING THE VALUES OF INTEGRALS

BY APPROXIMATION.

\medskip\rule{.1\linewidth}{.4pt}
\end{center}
\SeventhArticle{1}
Among the methods proposed for the approximate numerical determination of integrals, a prominent place is held by the rules developed by \textsc{Cotes}, following the lead of the great \textsc{Newton}. Namely, if the value of the integral $\int y\,\mathrm{d}x$ is required from $x=g$ to $x=h$, the values of $y$ at these extreme values of $x$ and at any number of intermediate values proceeding from the first to the last in equal increments must be multiplied by certain numerical coefficients. The sum of the products, multiplied by $h-g$, will then give the required integral, with greater accuracy the more terms are used in this operation. Since the principles of this method, which seems to be used by geometers less often than it deserves, have nowhere, as far as I know, been explained more fully, it will be consistent with our purpose to begin with a few remarks about them.

\SeventhArticle{2}
Let $n+1$ be the number of terms chosen for use, and let us set $h-g=\Delta$, so that the values of $x$ are $g$, $g+\frac{\Delta}{n}$, $g+\frac{2\Delta}{n}$, $g+\frac{3\Delta}{n}$, etc., up to $g+\Delta$, with the respective corresponding values of $y$ being $A$, $A'$, $A''$, $A'''$, etc., up to $A^{(n)}$. Finally, set $x=g+\Delta t$ in general, so that $y$ can also be regarded as a function of $t$. Let us denote by $Y$ the following function
\SourceLeaf{175}
\[
\begin{aligned}
&A.\;\frac{(nt-1)(nt-2)(nt-3)\ldotp\ldotp\ldotp\ldotp(nt-n)}{(-1)\quad(-2)\, .\,(-3)\ldotp\ldotp\ldotp\ldotp(-n)}\\
{}+{}&A'.\;\frac{nt\, .\,(nt-2)\, .\,(nt-3)\ldotp\ldotp\ldotp\ldotp\ldotp\ldotp(nt-n)}{1\, .\,(-1)\quad(-2)\ldotp\ldotp\ldotp\ldotp\ldotp\ldotp(1-n)}\\
{}+{}&A''.\;\frac{nt\, .\,(nt-1)\, .\,(nt-3)\ldotp\ldotp\ldotp\ldotp\ldotp\ldotp(nt-n)}{2.\quad1.\quad(-1)\ldotp\ldotp\ldotp\ldotp\ldotp\ldotp(2-n)}\\
{}+{}&A'''.\;\frac{nt\, .\,(nt-1)\, .\,(nt-2)\ldotp\ldotp\ldotp\ldotp\ldotp\ldotp(nt-n)}{3.\quad2.\quad1\ldotp\ldotp\ldotp\ldotp\ldotp\ldotp(3-n)}\\
{}+{}&\text{etc.}\\
{}+{}&A^{(n)}.\;\frac{nt(nt-1)\, .\,(nt-2)\ldotp\ldotp\ldotp\ldotp\ldotp\ldotp(nt-n+1)}{n\, .\,(n-1)\quad(n-2)\ldotp\ldotp\ldotp\ldotp\ldotp\ldotp1}
\end{aligned}
\]
\Needspace{10\baselineskip}or $\sum\dfrac{A^{(\mu)}T^{(\mu)}}{M^{(\mu)}}$, where $\mu$ represents each of the integers $0,1,2,3\ldotp\ldotp\ldotp n$,
\[
T^{(\mu)}=\frac{nt(nt-1)(nt-2)(nt-3)\ldotp\ldotp\ldotp\ldotp(nt-n)}{nt-\mu}
\]
\[
M^{(\mu)}\quad\text{value of }T\text{ for }nt=\mu.
\]
It is evident that $Y$ is a polynomial in $t$ of degree $n$, and that its values at each of the $n+1$ values of $t$, namely $0,\frac1n,\frac2n,\frac3n\ldotp\ldotp\ldotp\ldotp1$, are equal to the values of $y$. It is also clear that if $Y'$ is another polynomial agreeing with $y$ at the same values, $Y'-Y$ vanishes at those values and is therefore divisible by the factors $t,t-\frac1n,t-\frac2n,t-\frac3n\ldotp\ldotp\ldotp\ldotp t-1$, and consequently also by their product (which has degree $n+1$). Hence it is clear that $Y'$, unless it is entirely identical with $Y$, must certainly have a higher degree; in other words, $Y$ is the unique polynomial of degree not exceeding $n$ that, at those $n+1$ values, agrees with $y$. Thus, if $y$, when expanded in a series in powers of $t$, terminates entirely before the term involving $t^{n+1}$, it will be identical with $Y$. If, however, it at least converges so rapidly that the subsequent terms may be neglected, the function $Y$, between the limits $t=0,t=1$, or $x=g,x=h$, can take the place of $y$.

